\section{Shared Kernels}
\label{sec:kernels}

Diagrams do not each reinvent geometry. A small set of shared kernels under
\texttt{src/graph}, \texttt{src/text}, and \texttt{src/style} does the reusable
work; diagram layout engines compose them. This is what keeps the per-diagram
layout code short.

\subsection{Layout (\texttt{src/graph})}

\begin{description}
  \item[\texttt{layered.ts}] A Sugiyama-lite layered placement: longest-path
        ranking plus size-aware coordinate assignment. Used by the node-link and
        UML kinds (class, state, ER, flowchart) so they share rank assignment.
  \item[\texttt{tree.ts}] A tidy hierarchical placement: every parent is
        centred over the block its descendants occupy, sibling subtrees never
        overlap, and node sizes may vary. Used by the whole tree family.
  \item[\texttt{connect.ts}] Edge-endpoint helpers: \texttt{borderPoint} clips a
        connector to a node's border; \texttt{connectSlots} clips both ends along
        the line joining two boxes --- the \emph{slot-aware pointer} used for
        cell-to-cell, node-to-node, and panel-to-panel arrows alike.
\end{description}

\subsection{Text (\texttt{src/text})}

\texttt{measureText} gives deterministic width/height for a string at a font
size (layout must not depend on a live font engine), and \texttt{wrap}/truncate
helpers handle multi-line and overflow.

\subsection{Cells and strips (\texttt{src/scene/strip.ts})}

\texttt{buildStrip} produces a contiguous (or gapped) run of equal cells ---
arrays, heap array-backings, B-tree key strips, slotted pages --- and returns the
bounding rectangle of each cell as a \emph{slot} so pointers can clip to
individual cells.

\subsection{Cost tiers (\texttt{src/style/cost.ts})}

A reusable styling layer maps a numeric edge weight (latency, hop distance, plan
cost) onto a discrete visual tier --- colour plus dash pattern --- and builds a
legend:

\begin{lstlisting}[language=ts]
classifyCost(scale, weight): CostTier   // bucket a weight into a tier
buildLegend(pen, theme, scale, origin)  // render the tier legend
\end{lstlisting}

It is pure styling, no layout, and is consumed by the topology family
(\S\ref{sec:families}).
